\documentclass[11pt,letterpaper]{article}
\usepackage{technical-design-spec}
\setDocTitle{Deployment Promotion, Rollback and Recovery Architecture}
\setDocSubtitle{Architecture and security boundaries}
\setDocVersion{1.0}
\setDocStatus{Technical education / proposed reference}
\setDocOwner{CI/CD Engineering}
\setDocDate{30 October 2026}
\setDocSummary{Supplementary engineering practice}
\begin{document}
\makedoctitle
\makedocsummary
\tableofcontents
\clearpage
\section{Status and scope}
Proposed reference architecture. No production cloud environment or organizational authorization is implied. Illustrative diagrams reside in this suite and should be verified against the actual runner, GitHub and cloud trust configuration.

\section{Release identity}
Source revision, build workflow identity, toolchain lock, artifact cryptographic digest and approval record form release identity. Promote identical artifacts through environments; do not rebuild between staging and production unless signed policy explains why.

\section{Lifecycle}
Candidate package -> test evidence -> software composition/SAST review -> provenance/attestation verify -> staging deploy -> synthetic verification -> release approval -> production progressive deployment -> health observation and rollback decision.

\section{Database implications}
Code rollback may be unsafe after irreversible schema change. Prefer backward-compatible expand/migrate/contract sequences, audited backup and tested recovery window; require independent change authority for destructive schema operations.

\section{Protected execution}
Only approved protected branches/tags may request deploy; environment rules apply when supported; OIDC cloud trust policy must check the actual subject and audience. A deployment job retrieves immutable candidate by digest rather than taking a URL from untrusted PR content.

\section{Operational acceptance}
Define health endpoints and synthetic checks, error budget thresholds, observation windows, abort criteria, rollback owner, failure communication and audit capture before deploying.

\section{Acceptance review}
\begin{itemize}
\item Record context, container/module responsibilities and owning teams.
\item Identify untrusted inputs and every elevation or data transfer.
\item Exercise one successful path and at least one deliberate unauthorized path.
\item Identify release restoration point and hard-stop thresholds.
\item Attach reviewer sign-off and an architecture decision record for nontrivial trade-offs.
\end{itemize}
\end{document}
